\section{Discussion}

The expected policy implication is that domestic net zero should be treated as a distributional household risk-protection agenda. Protection depends on the alignment of fabric retrofit, heating technology, distributed generation, flexibility, tariff design, electricity market reform, capital access, and tenure-sensitive governance.

In Westminster, this means paying particular attention to private renters, leasehold flats, HMOs, heritage-constrained dwellings, district-heating governance, social housing delivery pathways, and households whose vulnerability is masked by high neighbourhood average wealth.
